% Generated from project outputs. Do not edit by hand.
\begin{table}[!htbp]
\centering
\begin{threeparttable}
\caption{Labor income by subsample: 2016-2019}
\label{tab:descriptive-2016-2019-subsamples-y1}
\small
\setlength{\tabcolsep}{4pt}
\renewcommand{\arraystretch}{1.15}
\begin{tabularx}{\linewidth}{@{}Xrrrrr@{}}
\toprule
Sample & N & Mean t & SD t & Mean t+1 & SD t+1 \\
\midrule
General & 276 & 2,607.3 & 1,835.7 & 2,685.5 & 1,792.5 \\
Men & 276 & 2,907.2 & 2,200.0 & 2,964.6 & 2,130.7 \\
Women & 276 & 2,093.2 & 1,454.8 & 2,213.9 & 1,530.2 \\
Urban & 276 & 2,883.1 & 1,869.6 & 2,949.1 & 1,811.5 \\
Rural & 271 & 1,922.4 & 1,420.2 & 2,085.0 & 1,570.0 \\
Indigenous & 271 & 1,900.5 & 1,314.9 & 1,997.3 & 1,351.3 \\
Non-Indigenous & 276 & 2,822.2 & 1,887.4 & 2,915.5 & 1,853.0 \\
Bottom 99 percent & 276 & 2,372.6 & 1,368.1 & 2,453.3 & 1,351.8 \\
Bottom 90 percent & 276 & 1,935.6 & 942.6 & 2,004.9 & 920.7 \\
Bottom 50 percent & 272 & 990.0 & 423.3 & 1,060.9 & 431.5 \\
\bottomrule
\end{tabularx}
\begin{tablenotes}[flushleft]
\footnotesize
\item Monthly income in quetzales. Unweighted descriptive statistics of survey-weighted cohort means. N counts cohort-transition rows, not individuals or independent cohorts. SD is dispersion across cohort means, not the standard error of a mean. Both incomes must be observed. Subsamples overlap; bottom-income groups are cumulative.
\end{tablenotes}
\end{threeparttable}
\end{table}
